% ============================================================================
% STATUS SECTION — rehearsal is a hypothesis. DelayProbeEnv is implemented, but
% no trained delay-probe checkpoint or causal delay-clamp result is reported.
% ============================================================================
\section{Attention-based rehearsal remains a hypothesis}
\label{sec:rehearsal}

A long-standing hypothesis holds that spatial working memory is maintained by
covertly shifting attention toward the memorized location --- the spatial analogue
of articulatory rehearsal --- so that the usual consequences of spatial attention
should be measurable at the remembered location during a delay
\cite{awh_jonides_reuterlorenz1998_rehearsal,awh_jonides2001_overlapping_attention_wm}.

The present model provides observations compatible with this hypothesis but not a
rehearsal result. It later re-acquires cue-directed attention after the blank, and
some checkpoints show elevated cued-location activity during a blank frame. No
cue-location decoder was fitted across the delay, so these are indirect signatures
compatible with maintenance rather than a decoded maintained-location result. In the
canonical affine VDA4 trace, however, the measured cued attention share falls near
the uniform baseline during the blank ($\alpha_c\approx0.22$ versus $0.25$). Neither
pattern by itself shows that attention is the load-bearing maintenance mechanism.

A delay-probe environment has been implemented, but we found no trained probe
checkpoint, memorized-location enhancement result, or causal delay-clamp analysis to
report. The decisive tests remain prospective: compare probe processing at cued and
uncued locations during the delay, then suppress cued-location attention during that
period and test whether any probe advantage or later detection benefit is reduced.
Until those measurements exist, attention-based rehearsal is a hypothesis rather
than an empirical claim of this manuscript.
